\documentclass[10pt,a4paper]{amsart}
\usepackage[margin=19mm]{geometry}
\usepackage{amsmath,amssymb,mathtools}
\usepackage[scr=rsfs,cal=euler]{mathalfa}
\usepackage{xfrac}
\usepackage[unicode,hidelinks,bookmarks=false]{hyperref}
\newcommand{\ie}{i.e.\ }
\newcommand{\cf}{cf.\ }
\newcommand{\resp}{respectively\ }
\newcommand{\Subsection}{\subsection}
\providecommand{\nobreakdash}{\nobreak\hbox{-}}
\setlength{\emergencystretch}{2em}
\multlinegap0pt
\usepackage{amssymb}
\usepackage{tikz}
\usepackage[shortlabels]{enumitem}
\usepackage{xcolor}
\newtheorem{theorem}{Theorem}
\title{Graph labellings and external difference families}
\author{Gavin Angus, Sophie Huczynska, Struan McCartney}
\date{Source: arXiv:2603.05662v1 (CC BY 4.0)}
\begin{document}
\maketitle

\noindent\emph{Source and licence.} Graph labellings and external difference families. Version arXiv:2603.05662v1 (CC BY 4.0), distributed by arXiv under the Creative Commons Attribution 4.0 International licence, http://creativecommons.org/licenses/by/4.0/, as recorded in the arXiv licence metadata for this exact version and shown on the abstract page https://arxiv.org/abs/2603.05662v1. Full licence notice: \texttt{licenses/arxiv-2603.05662-CC-BY-4.0.txt}.\par\smallskip

\noindent\textit{Editorial scope.} Counted unit: the article's theorem orblowsmall (source0.tex lines 952--954). The article's complete original proof of it is delivered with it (source0.tex lines 955--998, one \texttt{proof} environment of 44 lines). No mathematics is written, repaired or re-derived: the statement and the proof are the article's own lines, in the article's own notation, and no retained line is substituted or re-flowed.  No further source-local context is needed: every cross-reference of every retained line resolves inside the two delivered spans.  Nothing was omitted from any retained span, so the package carries no omission record.  No retained line contains a citation, so the wrapper carries no bibliography and no bibliography entry.  No retained line contains a figure environment, an image-inclusion command or drawing code, so no image is delivered and the package declares no asset dependency.  Editorial formatting only, in addition: an AMS \texttt{amsart} preamble that restates the article's own declarations and macros used by the retained lines, namely the environments \texttt{theorem} on separate counters, together with the standard packages those lines need.  The retained environments print in this document as Theorem 1 in order of appearance, whereas the article prints the same items with the numbering of its own full document.  The complete native source of the article as distributed in the arXiv e-print for this exact version is bundled unchanged as \texttt{sources/195816/source0.tex}.
\begin{theorem}\label{orblowsmall}
    Let $G$ be a digraph that has an oriented near $\alpha$-valuation $p$, and let $l$ be a positive integer. The blow-up $G_{l,1}$ has an oriented near $\alpha$-valuation.
\end{theorem}

\begin{proof}

Consider a directed edge $(u,v) \in \overrightarrow{E}(G)$; as $p$ is an oriented near $\alpha$-valuation, 
$p(v)-p(u) \equiv k \mod n+1$
for some $k\in \{1,\ldots,n\}$.  By the second condition for an oriented near $\alpha$-valuation, one of the vertices $u,v$ is in $V_{large}$ and the other is in $V_{small}$. For each $k \in \{1,\ldots,n\}$, it is obtained as the label $p(v)-p(u)$ of precisely one directed edge $(u,v)$ of $G$.  Denote by $P$ the set of $k\in \{1,\ldots,n\}$ which arise as positive integer differences (i.e. from $(u,v)$ where $u \in V_{small}$ and $v \in V_{large}$), and by $N$ the set of $k \in \{1,\ldots,n\}$ which arise as negative integer differences (i.e. from $(u,v)$ where $u \in V_{large}$ and $v \in V_{small}$).  We work in the integers until the final step of the proof.

The first case (i.e. labels from $P$) is exactly the same as in the near $\alpha$ case. 
\begin{itemize}
\item We replace $v$ with a singleton set of vertices labelled by $V_v=\{lp(v)\}$.
\item We replace $u$ with the $l$-set of vertices labelled by $V_u = \bigcup_{i=0}^{l-1}\{lp(u)+i\}$.
\item The directed edge $(u,v)$ is replaced with the $l$ directed edges from $V_u$ to $V_v$; the corresponding differences for these edges are:
$lp(v)i-lp(u)-i =  l(p(v)-p(u))-i = lk-i$
where $0 \leq i \leq l-1$ and $k \in P$. 
\end{itemize}

Taking multiset unions of all of these differences:
\[
\bigcup_{\substack{(u,v) \in G \\ u \in V_{small}}} \left( \bigcup_{i=0}^{l-1} \{ (l(p(v)-p(u))-i) \} \right)
= \bigcup_{k\in P} \bigcup_{i=0}^{l-1} \{ lk-i \} = \bigcup_{k \in P}[l(k-1)+1,lk]. 
\]

Now consider the directed edges with labels corresponding to $N$. 
\begin{itemize}
\item We replace $u$ with a singleton set of vertices labelled by $V_u=\{lp(u)\}$.
\item We replace $v$ with the $l$-set of vertices labelled by $V_v = \bigcup_{i=0}^{l-1}\{lp(v)+i\}$.
\item  The directed edge $(u,v)$ is replaced with the $l$ directed edges from $V_u$ to $V_v$.  Since $p(v)-p(u)$ is a negative integer, the actual integer difference corresponding to $k \in N$ in the $\mathbb{Z}_{n+1}$ setting was in fact $-(n+1-k)$.  So the corresponding integer differences for these edges in the blow-up setting are:
\[lp(v)+i-lp(u)=l(p(v)-p(u))+i = l(-(n+1-k))+i=lk-(l-1)+i-(nl+1).\]
where $0 \leq i \leq l-1$ and $k \in N$.  
\end{itemize}
Taking multiset unions of all of these differences:
\[
\bigcup_{(u,v) \in G, v \in V_{small}} \left( \bigcup_{i=0}^{l-1} \{ (l(p(v)-p(u)-(n-1))+i) \} \right)
= \bigcup_{k\in N} \bigcup_{i=0}^{l-1} \{ lk-(l-1)+i-(nl+1) \}\]
\[= \bigcup_{k \in N}[l(k-1)+1-(nl+1),lk-(nl+1)] \equiv \bigcup_{k \in N}[l(k-1)+1,lk]\mod nl+1. \]

It is clear that all edges in the blow-up are formed by one of the two above cases. Taking the multiset union over $k \in P$ and $k \in N$ gives all $k \in \{1, \ldots ,n \}$, so (working now in $\mathbb{Z}_{nl+1}$):
\[\bigcup_{(x,y) \in G_{l,1}} \{p'(y)-p'(x)\} \equiv \bigcup_{k=1}^n[l(k-1)+1,lk] \mod nl+1 = \mathbb{Z}_{nl+1} \setminus \{0\}.\]

Now note that, since $p$ is an oriented near $\alpha$-valuation, for each $ y \in V_{large}$ and $x \in V_{small}$ we have $p(y) > p(x)$ (in the ordering $0<1< \cdots < n$). Hence in $G_{l,1}$ for any $s \in V_y$ and $t \in V_x$:
    \[p'(s) \geq p(y)l >   p(x)l = p'(t)\]
(in the ordering $0<1< \cdots< nl$).
Hence $p'$ is an oriented near $\alpha$-valuation.

\end{proof}

\end{document}
